% Options for packages loaded elsewhere
\PassOptionsToPackage{unicode}{hyperref}
\PassOptionsToPackage{hyphens}{url}
%
\documentclass[
  11pt,
]{article}
\usepackage{amsmath,amssymb}
\usepackage{iftex}
\ifPDFTeX
  \usepackage[T1]{fontenc}
  \usepackage[utf8]{inputenc}
  \usepackage{textcomp} % provide euro and other symbols
\else % if luatex or xetex
  \usepackage{unicode-math} % this also loads fontspec
  \defaultfontfeatures{Scale=MatchLowercase}
  \defaultfontfeatures[\rmfamily]{Ligatures=TeX,Scale=1}
\fi
\usepackage{lmodern}
\ifPDFTeX\else
  % xetex/luatex font selection
\fi
% Use upquote if available, for straight quotes in verbatim environments
\IfFileExists{upquote.sty}{\usepackage{upquote}}{}
\IfFileExists{microtype.sty}{% use microtype if available
  \usepackage[]{microtype}
  \UseMicrotypeSet[protrusion]{basicmath} % disable protrusion for tt fonts
}{}
\makeatletter
\@ifundefined{KOMAClassName}{% if non-KOMA class
  \IfFileExists{parskip.sty}{%
    \usepackage{parskip}
  }{% else
    \setlength{\parindent}{0pt}
    \setlength{\parskip}{6pt plus 2pt minus 1pt}}
}{% if KOMA class
  \KOMAoptions{parskip=half}}
\makeatother
\usepackage{xcolor}
\usepackage[margin=1in]{geometry}
\usepackage{longtable,booktabs,array}
\usepackage{calc} % for calculating minipage widths
% Correct order of tables after \paragraph or \subparagraph
\usepackage{etoolbox}
\makeatletter
\patchcmd\longtable{\par}{\if@noskipsec\mbox{}\fi\par}{}{}
\makeatother
% Allow footnotes in longtable head/foot
\IfFileExists{footnotehyper.sty}{\usepackage{footnotehyper}}{\usepackage{footnote}}
\makesavenoteenv{longtable}
\setlength{\emergencystretch}{3em} % prevent overfull lines
\providecommand{\tightlist}{%
  \setlength{\itemsep}{0pt}\setlength{\parskip}{0pt}}
\setcounter{secnumdepth}{-\maxdimen} % remove section numbering
% definitions for citeproc citations
\NewDocumentCommand\citeproctext{}{}
\NewDocumentCommand\citeproc{mm}{%
  \begingroup\def\citeproctext{#2}\cite{#1}\endgroup}
\makeatletter
 % allow citations to break across lines
 \let\@cite@ofmt\@firstofone
 % avoid brackets around text for \cite:
 \def\@biblabel#1{}
 \def\@cite#1#2{{#1\if@tempswa , #2\fi}}
\makeatother
\newlength{\cslhangindent}
\setlength{\cslhangindent}{1.5em}
\newlength{\csllabelwidth}
\setlength{\csllabelwidth}{3em}
\newenvironment{CSLReferences}[2] % #1 hanging-indent, #2 entry-spacing
 {\begin{list}{}{%
  \setlength{\itemindent}{0pt}
  \setlength{\leftmargin}{0pt}
  \setlength{\parsep}{0pt}
  % turn on hanging indent if param 1 is 1
  \ifodd #1
   \setlength{\leftmargin}{\cslhangindent}
   \setlength{\itemindent}{-1\cslhangindent}
  \fi
  % set entry spacing
  \setlength{\itemsep}{#2\baselineskip}}}
 {\end{list}}
\usepackage{calc}
\newcommand{\CSLBlock}[1]{\hfill\break\parbox[t]{\linewidth}{\strut\ignorespaces#1\strut}}
\newcommand{\CSLLeftMargin}[1]{\parbox[t]{\csllabelwidth}{\strut#1\strut}}
\newcommand{\CSLRightInline}[1]{\parbox[t]{\linewidth - \csllabelwidth}{\strut#1\strut}}
\newcommand{\CSLIndent}[1]{\hspace{\cslhangindent}#1}
\ifLuaTeX
  \usepackage{selnolig}  % disable illegal ligatures
\fi
\usepackage{bookmark}
\IfFileExists{xurl.sty}{\usepackage{xurl}}{} % add URL line breaks if available
\urlstyle{same}
\hypersetup{
  pdftitle={CT-RAG: Structured Experiential Memory via Causal-Topological Retrieval for Long-Lived Agents},
  pdfauthor={Jean Carlo Nascimento},
  hidelinks,
  pdfcreator={LaTeX via pandoc}}

\title{CT-RAG: Structured Experiential Memory via Causal-Topological
Retrieval for Long-Lived Agents}
\author{Jean Carlo Nascimento}
\date{September 2026}

\begin{document}
\maketitle

\section{Candidate titles}\label{candidate-titles}

\begin{enumerate}
\def\labelenumi{\arabic{enumi}.}
\tightlist
\item
  \textbf{CT-RAG: Structured Experiential Memory via Causal-Topological
  Retrieval for Long-Lived Agents} \emph{(selected)}
\item
  \textbf{Beyond Similarity and Recency: Runtime Causal Provenance for
  Long-Lived Agent Memory}
\item
  \textbf{Retrieval Becomes Navigation: Causal-Topological Memory and
  Dynamic Terrain for Agents}
\end{enumerate}

\section{Abstract}\label{abstract}

Long-lived software agents accumulate execution histories that are
difficult to use as memory. Conventional retrieval-augmented generation
(RAG) primarily ranks items by semantic similarity, while temporal
memory systems emphasize when facts were valid or observed. Neither
geometry is sufficient when the relevant evidence is defined by runtime
execution structure: an event can be causally authoritative despite
being semantically dissimilar, temporally distant, or received out of
order. We introduce \textbf{CT-RAG}, a retrieval architecture for
\textbf{structured experiential memory} that represents execution
evidence using typed causal, temporal, and behavioral relations,
separates authoritative causal provenance from inferred association, and
maintains a dynamic, non-authoritative \textbf{terrain} whose
navigational influence can be reinforced or eroded without rewriting
immutable historical evidence. Retrieval is therefore treated as
navigation over preserved experience, while learning changes navigation
rather than truth.

We evaluate CT-RAG using controlled adversarial fixtures designed to
separate causal structure from simpler ranking signals. In an
out-of-order causation experiment, a child event is ingested before its
parent, with a 336-hour event-time gap and 25 semantically adversarial
distractors. Semantic-only and recency-only retrieval rank the true
parent 26th of 26 candidates, while CT-RAG \texttt{WHY} ranks it first using explicit
runtime \texttt{causation\_id} provenance. In an eroded-path recovery
experiment, a raw-success baseline prefers an unreachable
\texttt{GlobalFix} path with 9/10 historical success, whereas CT-RAG
excludes it because it is not causally reachable from the current state
and ranks the best reachable recovery first. Conservative small-sample
support is represented by the 95\% Wilson lower bound rather than an ad
hoc count bonus. CT-RAG achieves Precision@1 = 1.0 in this controlled
recovery fixture, while terrain-only, recency-only, and raw-success
baselines each obtain 0.0; false-rescue increases from 0.0 at k=1 to
0.75 at k=4, exposing rather than hiding degradation as more
alternatives are returned. A secondary same-logic degradation study under the current zero-FPR null
yields a median +1 day relative lead across a 28-cell parameter surface,
with the positive-lead cells forming one contiguous region and the
+1-day lead invariant to four phase offsets. These experiments are
mechanism-level synthetic evidence, not production generalization. The
results support a narrower claim: \textbf{runtime causal reachability
contains retrieval information that semantic similarity, recency, and
global success statistics do not encode}.

\section{1. Introduction}\label{introduction}

Retrieval-augmented generation augments a parametric language model with
information retrieved from an external memory
(\citeproc{ref-lewis2020rag}{Lewis et al. 2020}). Most deployed and
research RAG systems begin from a geometric assumption: if two items are
close in an embedding space, one is likely to be useful for answering a
query. Long-term agent-memory systems add persistence, reflection,
hierarchical storage, temporal knowledge graphs, or graph navigation
(\citeproc{ref-park2023generative}{Park et al. 2023};
\citeproc{ref-packer2023memgpt}{Packer et al. 2023};
\citeproc{ref-rasmussen2025zep}{Rasmussen et al. 2025};
\citeproc{ref-gutierrez2024hipporag}{Gutiérrez et al. 2024}). This
progression is important, but agents operating inside event-driven
systems face an additional problem. Their relevant memory is often not
merely \emph{about the same thing} or \emph{nearby in time}. It is the
execution that caused the current event, the previously successful route
reachable from the current state, the basin toward which repeated
behavior is drifting, or the historical branch whose current
navigational influence has eroded without invalidating the evidence that
it once occurred.

Distributed systems make the distinction concrete. Event arrival order
is not event time, event time is not causal order, and semantic
similarity is not causal authority. Lamport's happens-before relation
established the importance of partial order over physical-clock
proximity in distributed computation
(\citeproc{ref-lamport1978time}{Lamport 1978}). Stream-processing
systems similarly distinguish event time from processing time because
out-of-order arrival is intrinsic rather than exceptional
(\citeproc{ref-akidau2015dataflow}{Akidau et al. 2015}). Temporal
databases distinguish valid time from transaction time for related
reasons (\citeproc{ref-jensen2018temporal}{Jensen and Snodgrass 2018}).
CT-RAG adopts the same discipline for agent memory: \textbf{different
relations carry different semantics and must not silently collapse into
one another}.

This paper is organized around two distinct functions of experiential
topology. \textbf{RQ1 treats topology as a relation of authority:} when
several memories are semantically or temporally plausible, which prior
event is entitled to explain the current event because the runtime
explicitly identifies it as causal provenance? \textbf{RQ2 treats
topology as a relation of feasibility:} among historically successful
recovery routes, which alternatives are actually reachable from the
current state? Similarity, recency, and global success statistics can
rank observations, but they do not encode either authority or
state-conditioned feasibility. The two primary experiments are designed
to isolate these two roles separately.

We use the term \textbf{structured experiential memory} to denote memory
organized around observed execution experience rather than only facts or
text fragments. This term is intentionally descriptive rather than a
claim that graph-structured experience is unique to CT-RAG. Recent work
explicitly studies graph-based agent memory and experiential memory
(\citeproc{ref-yang2026graphmemory}{Yang et al. 2026};
\citeproc{ref-hu2025memorysurvey}{Hu et al. 2025};
\citeproc{ref-dai2026gsem}{Dai et al. 2026}), while other 2026 systems
explore retrieval-driven reconsolidation and selective forgetting
(\citeproc{ref-song2026realm}{Song et al. 2026};
\citeproc{ref-rusu2026selective}{Rusu, Khanzadeh, and Alalfi 2026}). Our
claim is narrower: CT-RAG combines runtime-owned causal provenance,
typed separation of experiential relations, execution basins/attractors,
and a dynamic terrain overlay that changes navigational influence
without changing historical truth.

The design follows one central invariant:

\begin{quote}
\textbf{Observed adjacency may create temporal or behavioral evidence,
but only explicit or separately justified provenance may create a causal
edge.}
\end{quote}

This distinction also separates CT-RAG from causal inference. A runtime
field such as \texttt{causation\_id} is not a statistical estimate that
one variable causes another; it is provenance asserted by the executing
system. Conversely, observational alternatives retrieved from history do
not identify intervention effects in Pearl's sense
(\citeproc{ref-pearl2009causality}{Pearl 2009}). Granger-style
predictability and PCMCI-style causal discovery can be useful analytical
tools, but they answer different questions and are not substituted for
runtime provenance (\citeproc{ref-granger1969causal}{Granger 1969};
\citeproc{ref-runge2020pcmci}{Runge 2020}).

This paper makes four contributions.

\begin{enumerate}
\def\labelenumi{\arabic{enumi}.}
\tightlist
\item
  \textbf{A typed experiential topology.} We model agent/system
  experience as nodes connected by causal, temporal, behavioral, and
  optionally semantic relations, with causal edges carrying explicit
  provenance and temporal relations carrying scope.
\item
  \textbf{A non-authoritative dynamic terrain.} Repeated trajectories
  reinforce navigational influence and unused paths erode, while the
  underlying nodes, edges, provenance, and observed transition history
  remain intact.
\item
  \textbf{Mode-aware causal-topological retrieval.} \texttt{WHY}
  retrieval can prioritize authoritative causal ancestry even under
  adversarial semantic and temporal geometry; \texttt{RECOVERY} uses
  causal reachability as an eligibility gate and conservative historical
  success support to order reachable alternatives without letting low
  terrain influence act as a veto.
\item
  \textbf{Adversarial controlled evidence.} We construct baselines
  intended to defeat the method: 25 semantically near and temporally
  recent distractors in the causal-gap experiment, and a globally
  higher-success but causally unreachable recovery in the eroded-path
  experiment. These fixtures isolate information carried by runtime
  topology rather than merely showing that an implementation satisfies
  its own specification.
\end{enumerate}

The remainder of the paper describes the model, implementation
invariants, experimental protocol, results, and limitations. All
numerical results reported here were regenerated by GitHub Actions CI
run \textbf{\#210} from commit
\textbf{\texttt{d2c73de6ae1a691ad7d61c9ece26a4e84c93c44c}} on Python
3.11, 3.12, and 3.13; the tables below use the Python 3.13 uploaded
benchmark artifact.

\section{2. Related Work}\label{related-work}

\subsection{2.1 Retrieval-augmented generation and graph
retrieval}\label{retrieval-augmented-generation-and-graph-retrieval}

RAG combines parametric generation with non-parametric retrieval
(\citeproc{ref-lewis2020rag}{Lewis et al. 2020}). Dense retrieval is
effective when relevance correlates with semantic similarity, but it
does not by itself encode execution causality, temporal validity, or
reachability constraints.

HippoRAG combines LLM extraction, a knowledge graph, and Personalized
PageRank to support efficient multi-hop retrieval and long-term
knowledge integration (\citeproc{ref-gutierrez2024hipporag}{Gutiérrez et
al. 2024}). Its contribution demonstrates that graph structure can
improve retrieval beyond flat vector search. CT-RAG does not claim
generic graph navigation as novel. The distinction is the graph's
semantics: CT-RAG's primary structure is \textbf{runtime experience},
and its causal relations are admitted under provenance constraints
rather than created because entities co-occur or are semantically
associated.

Topo-RAG argues that hybrid enterprise documents should preserve
intrinsic data topology rather than linearizing every structure into
text; it routes narrative and tabular evidence through
structure-appropriate retrieval
(\citeproc{ref-dantart2026toporag}{Dantart and Kóvacs-Navarro 2026}).
CT-RAG uses \emph{topology} in a different sense. Topo-RAG preserves the
intrinsic spatial/structural topology of an artifact; CT-RAG models the
topology of \textbf{execution experience} across events and states. An
important invariant follows: intrinsic artifact adjacency, temporal
adjacency, behavioral continuity, and causal evidence are separate
relations.

\subsection{2.2 Long-term agent memory}\label{long-term-agent-memory}

Generative Agents stores a stream of natural-language experiences,
retrieves memories dynamically, and synthesizes reflections that
influence future behavior (\citeproc{ref-park2023generative}{Park et al.
2023}). MemGPT manages multiple memory tiers using an operating-system
analogy to extend effective context over long interactions
(\citeproc{ref-packer2023memgpt}{Packer et al. 2023}). Zep/Graphiti uses
a temporally aware knowledge graph to dynamically integrate
conversational and business data while preserving historical
relationships (\citeproc{ref-rasmussen2025zep}{Rasmussen et al. 2025}).
These systems establish long-term memory as a first-class agent
primitive and motivate explicit treatment of memory dynamics.

Recent surveys distinguish factual, experiential, working, structural,
and graph-based memory, showing that graph structure is now a broad
design family rather than a differentiating claim by itself
(\citeproc{ref-hu2025memorysurvey}{Hu et al. 2025};
\citeproc{ref-yang2026graphmemory}{Yang et al. 2026}). GSEM, for
example, encodes multi-agent coordination episodes as heterogeneous
relational graphs and retrieves structurally similar experience to
accelerate adaptation (\citeproc{ref-dai2026gsem}{Dai et al. 2026}).
REALM reorganizes a heterogeneous memory graph based on retrieval
feedback (\citeproc{ref-song2026realm}{Song et al. 2026}). Selective
Forgetting reports a useful negative result: a graph-based
conversational memory pipeline did not automatically beat a flat vector
baseline under matched retrieval budget, while selective pruning could
reduce memory with little measured loss
(\citeproc{ref-rusu2026selective}{Rusu, Khanzadeh, and Alalfi 2026}).
These findings reinforce our claim discipline: \textbf{CT-RAG does not
assert that graphs are intrinsically better than vectors.} Its
experiments instead ask whether specific runtime relations carry
information absent from semantic or temporal baselines.

\subsection{2.3 Time and causality in distributed
systems}\label{time-and-causality-in-distributed-systems}

Lamport formalized happens-before as a partial order independent of
physical clock equality (\citeproc{ref-lamport1978time}{Lamport 1978}).
Modern dataflow systems distinguish event time, processing time, and
completeness/watermarks to handle unbounded out-of-order streams
(\citeproc{ref-akidau2015dataflow}{Akidau et al. 2015}). Temporal
database research similarly distinguishes valid time and transaction
time (\citeproc{ref-jensen2018temporal}{Jensen and Snodgrass 2018}).
CT-RAG draws an engineering lesson from this literature: memory systems
should not treat one scalar timestamp as a universal relation.

CT-RAG further distinguishes \textbf{runtime causal provenance} from
\textbf{causal inference}. Pearl's structural causal models and
intervention semantics concern identification of causal effects
(\citeproc{ref-pearl2009causality}{Pearl 2009}). Granger causality
concerns predictive temporal precedence
(\citeproc{ref-granger1969causal}{Granger 1969}), while PCMCI/PCMCI+
estimates lagged and contemporaneous causal relations from observational
time series under explicit assumptions
(\citeproc{ref-runge2020pcmci}{Runge 2020}). CT-RAG may store hypotheses
from such methods with lower epistemic status, but an event relation
asserted by \texttt{causation\_id} is a different object: it is
provenance emitted by the runtime, not an observational causal estimate.

\subsection{2.4 Memory strength as a limited
analogy}\label{memory-strength-as-a-limited-analogy}

Bjork and Bjork's New Theory of Disuse distinguishes storage strength
from retrieval strength (\citeproc{ref-bjork1992disuse}{Bjork and Bjork
1992}). CT-RAG's separation between preserved history and current
navigational influence is analogous in a limited engineering sense: a
path can become hard to retrieve without being erased. We do
\textbf{not} claim cognitive or psychological isomorphism. The terrain
is an engineered retrieval overlay, not a model of human memory.

\section{3. Problem Definition}\label{problem-definition}

Let an experiential memory be a typed directed graph

{[} G = (V, E\_c, E\_t, E\_b, E\_s), {]}

where (V) contains observed events or states; (E\_c) contains causal
relations; (E\_t) temporal relations; (E\_b)
behavioral/execution-continuity relations; and (E\_s) optional semantic
relations. Semantic and lexical similarity may also be computed as
retrieval channels without materializing (E\_s).

A causal edge is represented as

{[} e\_c = (u,v,p,q,\mathcal{E}), {]}

where (p) is a provenance class, (q\in[0,1]) a confidence, and
(\mathcal{E}) evidence pointers. The implementation distinguishes
provenance such as event, execution, workflow, dependency, inferred, and
hypothesized. The epistemic invariant is that an edge in (E\_c) cannot
exist without causal provenance.

A temporal edge is

{[} e\_t=(u,v,\sigma), {]}

where (\sigma) is a temporal scope. Implemented scopes are
\texttt{EXECUTION}, \texttt{ACTOR\_OR\_AGGREGATE}, \texttt{DEPLOYMENT},
\texttt{INCIDENT\_WINDOW}, and \texttt{GLOBAL\_OBSERVED}. The fail-safe
default for temporal traversal is execution-local; callers must
explicitly opt into broader scope traversal. This prevents accidental
chaining of unrelated executions merely because timestamps are adjacent.

For a query (q), anchor set (A), and mode (m), the base retriever
combines semantic, lexical, causal, topological, temporal, and
behavioral components:

{[} S\_\{base\}(v\mid q,A,m)= \sum\_j w\^{}\{(m)\}\_j s\_j(v\mid q,A),
{]}

where the weight vector depends on query mode. CT-RAG currently exposes
\texttt{SIMILAR}, \texttt{WHY}, \texttt{WHAT\_NEXT}, \texttt{RECOVERY},
and \texttt{COUNTERFACTUAL}. \texttt{COUNTERFACTUAL} is observational
retrieval of historical divergence and must not be interpreted as an
identified intervention effect.

\section{4. Architecture and
Invariants}\label{architecture-and-invariants}

\subsection{4.1 Event projection and authoritative
causation}\label{event-projection-and-authoritative-causation}

The \texttt{EventProjector} maps authoritative event records into memory
nodes. When event (v) declares \texttt{causation\_id\ =\ u} and (u)
already exists, the projector creates

{[} u \xrightarrow{CAUSAL,\ EVENT} v {]}

with evidence source \texttt{event.causation\_id}. If (u) has not yet
arrived, the child is placed in a pending-causation structure keyed by
the parent identifier. When (u) later arrives, the edge is reconciled
deterministically. Arrival order therefore affects when the edge can be
materialized, but not its eventual causal authority.

Events sharing an execution identifier may also receive temporal and
behavioral edges according to execution sequence. Crucially, sequence
alone never creates a causal edge. This prevents the common error of
converting ``B happened after A'' into ``A caused B.''

\subsection{4.2 Degraded mode without causal
provenance}\label{degraded-mode-without-causal-provenance}

External telemetry often lacks authoritative causal fields. CT-RAG
therefore operates in a degraded but explicit mode: trace hierarchy,
temporal adjacency, and behavioral continuity may be represented, but
disconnected spans or events are not promoted into (E\_c) merely because
no competing explanation exists. This fail-safe behavior reduces recall
of causal paths when provenance is missing, but preserves epistemic
integrity.

\subsection{4.3 Basins and attractors}\label{basins-and-attractors}

CT-RAG models repeated execution as movement through topology. A
structural sink is a node with no outgoing basin-relevant edge. A
recurrent strongly connected component (SCC) may also act as an
attractor. Basins are defined as bounded incoming neighborhoods of
registered attractors over causal and behavioral relations.

Given two snapshots with basin member sets (B\_a\^{}\{(0)\}) and
(B\_a\^{}\{(1)\}) for attractor (a), basin drift is represented by
Jaccard distance:

{[}
D\_B(a)=1-\frac{|B_a^{(0)}\cap B_a^{(1)}|}{|B_a^{(0)}\cup B_a^{(1)}|}.
{]}

This is a descriptive topological drift measure, not a causal-effect
estimate.

\section{5. Dynamic Terrain}\label{dynamic-terrain}

The graph records historical structure; the \textbf{DynamicTerrain}
records current navigational influence. For each edge (e), terrain
maintains an observation count (n\_e) and influence (I\_e). Observing a
transition reinforces only the overlay:

{[} I\_e \leftarrow \min(I\_\{max\}, I\_e + \alpha). {]}

After elapsed time (\Delta t), influence decays exponentially:

{[} I\_e(t+\Delta t)=\max(I\_\{min,e\}, I\_e(t)e\^{}\{-\mu\Delta t\}).
{]}

Protected rare-but-critical edges may use a non-zero floor
(I\_\{min,e\}). Decay never removes the underlying node, edge, causal
evidence, or transition count.

For a causal path (\pi=(e\_1,\ldots,e\_k)), path influence is the
geometric mean

{[} I(\pi)=\left(\prod\emph{\{i=1\}\^{}\{k\}
I}\{e\_i\}\right)\^{}\{1/k\}. {]}

Outside recovery mode, terrain can act as a multiplicative navigational
reranker. In \texttt{RECOVERY}, however, low influence is deliberately
\textbf{not} a veto: an old but historically successful reachable route
may need to be surfaced precisely because it has fallen out of current
use.

The implementation also supports a transition-surprise signal using a
Laplace-smoothed empirical outgoing probability,

{[} \mathrm{surprise}(e)=1-P(e\mid source(e),kind(e)), {]}

but surprise-based reinforcement is not required for the primary results
in this paper.

\section{6. Recovery Ranking: Reachability and Conservative Historical
Support}\label{recovery-ranking-reachability-and-conservative-historical-support}

Recovery retrieval answers a different question from similarity search:
\emph{from the current state, which historically observed routes can
reach a desired recovered state, and which of those routes have credible
evidence of success?}

CT-RAG first requires causal reachability from the current anchor. A
successful historical terminal with no causal path from the current
state is ineligible, regardless of global success rate. For a reachable
terminal, the final branch immediately preceding that terminal provides
counts (s) successful outcomes among (n) observed outcomes.

Rather than rank by the empirical ratio (\hat p=s/n) alone, CT-RAG uses
the lower endpoint of the 95\% Wilson score interval:

{[} LB\_W(s,n)=
\frac{\hat p+\frac{z^2}{2n}-z\sqrt{\frac{\hat p(1-\hat p)}{n}+\frac{z^2}{4n^2}}}
\{1+\frac{z^2}{n}\},\qquad z=1.96. {]}

This folds small-sample uncertainty into the ranking without inventing
an arbitrary support multiplier. For example, (LB\_W(1,1)\approx0.207),
while (LB\_W(8,8)\approx0.676); both empirical rates are 1.0, but the
latter has substantially stronger support.

The terrain-aware recovery reranker retains \texttt{terrain\_influence}
for explanation but adds conservative historical evidence only to
causally reachable successful terminals. Conceptually,

{[} S\_\{REC\}(v)=S\_\{base\}(v)+\lambda,LB\_W(s\_v,n\_v), {]}

subject to

{[} Reachable(A,v)=1. {]}

The current implementation uses (\lambda=0.35). This coefficient is an
implementation parameter, not claimed as universally optimal.

\section{7. Experimental Method}\label{experimental-method}

We separate \textbf{comparative experiments} from
\textbf{property/mechanism tests}. Comparative experiments include
competing conditions intended to succeed under plausible alternative
assumptions. Property tests verify invariants and regression behavior
but are not presented as evidence of superiority.

All primary results below are synthetic controlled fixtures. They are
designed for causal identification of \emph{mechanism behavior} inside
the software artifact, not for estimating production effect sizes.

\subsection{7.1 RQ1: Can runtime causal provenance recover evidence that
semantic and recency geometry
suppress?}\label{rq1-can-runtime-causal-provenance-recover-evidence-that-semantic-and-recency-geometry-suppress}

We ingest child event \texttt{Action\_B\_Success} first with explicit

\begin{verbatim}
causation_id = Action_A_Intent
\end{verbatim}

while the parent is absent. We then ingest 25 unrelated events whose
type is deliberately adversarial to semantic retrieval:
\texttt{Action.B.Success.Context}. The parent finally arrives with event
time 336 hours earlier than the child.

We compare:

\begin{itemize}
\tightlist
\item
  \textbf{Semantic-only:} global dense retrieval over all nodes.
\item
  \textbf{Recency-only:} descending event-time ordering.
\item
  \textbf{CT-RAG WHY:} mode-aware retrieval anchored at
  \texttt{Action\_B\_Success}, with causal provenance available from the
  reconciled \texttt{causation\_id} edge.
\end{itemize}

The primary metric is absolute rank of the true parent. The test oracle
requires CT-RAG parent rank = 1 and both baselines to place the parent
outside the top 10.

\subsection{7.2 RQ2: Does causal reachability add information beyond
terrain, recency, and raw historical
success?}\label{rq2-does-causal-reachability-add-information-beyond-terrain-recency-and-raw-historical-success}

The current state is \texttt{retry}. Reachable historical recovery
branches are:

\begin{itemize}
\tightlist
\item
  \texttt{ProviderFallback}: 8/10 successful outcomes;
\item
  \texttt{ManualPatch}: 1/10;
\item
  \texttt{ScriptPatch}: 2/10;
\item
  \texttt{CacheReset}: 1/4.
\end{itemize}

A separate \texttt{GlobalFix} branch has 9/10 historical success but
originates from a different root and is \textbf{not causally reachable
from \texttt{retry}}.

All historical recovery branches are eroded, while the current terrain
is dominated by \texttt{Timeout\ -\textgreater{}\ HumanIntervention}.

We compare:

\begin{itemize}
\tightlist
\item
  \textbf{Terrain-only:} current path influence;
\item
  \textbf{Recency-only:} most recent recovered terminal first;
\item
  \textbf{Raw-success-global:} global empirical (P(success)), ignoring
  reachability;
\item
  \textbf{Structural reachable ceiling:} all recovered terminals
  causally reachable from \texttt{retry}, deliberately treated as an
  unranked candidate set;
\item
  \textbf{CT-RAG RECOVERY:} reachability gate plus Wilson-supported
  historical success.
\end{itemize}

The gold top-1 recovery is \texttt{ProviderFallback}. We report
Precision@1 and false-rescue@k for (k\in\{1,2,3,4\}). The broader
false-rescue curve matters because a healer consuming multiple
suggestions is exposed to harmful alternatives even when top-1 is
correct.

\subsection{7.3 RQ3: Is early behavioral degradation detectable before
an FPR-matched infrastructure
detector?}\label{rq3-is-early-behavioral-degradation-detectable-before-an-fpr-matched-infrastructure-detector}

This secondary experiment uses daily windows whose terminal absorption
distribution shifts gradually from \texttt{Recovered} toward
\texttt{HumanIntervention}. Behavioral drift is total variation from a
fixed healthy baseline. To avoid the original unfair comparison against
a frozen infrastructure threshold, both behavioral and infrastructure
signals use the same sustained-change detection pattern, and the
infrastructure threshold is calibrated on stationary null scenarios to
match the behavioral false-positive rate for each parameter cell.

The robustness grid contains 28 combinations of behavioral drift
threshold and sustained windows. We report the contiguous positive-lead
region, lead distribution, FPR gap, and phase-offset sensitivity. The
phase analysis shifts window boundaries by fractions 0, 0.25, 0.50, and
0.75. This is an alignment test, not sub-daily evidence.

\subsection{7.4 Mechanism/property
validation}\label{mechanismproperty-validation}

Additional CI-gated tests validate that:

\begin{itemize}
\tightlist
\item
  decay changes influence but preserves graph history;
\item
  newly observed recurrent cycles increase SCC internal observations and
  confidence;
\item
  out-of-order causal reconciliation is deterministic;
\item
  execution-local temporal traversal is the fail-safe default.
\end{itemize}

These support implementation correctness but are not counted as
comparative experimental wins.

\section{8. Results}\label{results}

\subsection{8.1 Out-of-order causal
provenance}\label{out-of-order-causal-provenance}

\textbf{Table 1. Absolute rank of the true causal parent among 26 total
candidates under adversarial semantic and temporal geometry.}

\begin{longtable}[]{@{}lrrl@{}}
\toprule\noalign{}
Method & True parent rank & Top-10 recall & Information used \\
\midrule\noalign{}
\endhead
\bottomrule\noalign{}
\endlastfoot
Semantic-only dense retrieval & 26 & 0 & semantic similarity \\
Recency-only & 26 & 0 & event-time recency \\
CT-RAG \texttt{WHY} & \textbf{1} & \textbf{1} & runtime causal
provenance + topology \\
\end{longtable}

The reconciled edge carries provenance \texttt{EVENT}, evidence source
\texttt{event.causation\_id}, and causal component 0.9 in the retrieved
parent hit. The 336-hour temporal gap and 25 semantically close
distractors therefore do not prevent the parent from ranking first.

This experiment establishes a specific claim rather than generic
retrieval superiority. Semantic similarity and recency fail because the
fixture is constructed so that both geometries point toward distractors.
CT-RAG succeeds because the runtime supplies a relation those baselines
do not represent. The result is therefore evidence that \textbf{runtime
causal provenance can carry retrieval information orthogonal to semantic
and temporal proximity}.

\subsubsection{First-page figure}\label{first-page-figure}

We propose Figure 1 as a three-panel visual occupying the top half of
the first page after the abstract.

\begin{itemize}
\tightlist
\item
  \textbf{Panel A --- Semantic geometry:} \texttt{Action\_B\_Success} at
  the center of a dense cluster of 25 \texttt{Action.B.Success.Context}
  distractors; \texttt{Action\_A\_Intent} visually distant. Label:
  \emph{semantic-only parent rank = 26}.
\item
  \textbf{Panel B --- Recency geometry:} horizontal event-time axis
  showing \texttt{Action\_A\_Intent} 336 h to the left,
  \texttt{Action\_B\_Success}, then the 25 recent distractors. Label:
  \emph{recency-only parent rank = 26}.
\item
  \textbf{Panel C --- Runtime provenance:} collapse the same
  observations into a causal view with a single highlighted
  \texttt{Action\_A\_Intent\ -\/-causation\_id-\/-\textgreater{}\ Action\_B\_Success}
  edge crossing the temporal gap. Label: \emph{CT-RAG WHY parent rank =
  1}.
\end{itemize}

The caption should state: \textbf{Same evidence, three geometries.
Similarity and recency concentrate on adversarial distractors; explicit
runtime provenance identifies the causal parent.}

\subsection{8.2 Eroded-path rescue}\label{eroded-path-rescue}

The global raw-success baseline ranks \texttt{GlobalFix} first because
its empirical success is 0.90, above \texttt{ProviderFallback} at 0.80.
However, \texttt{GlobalFix} is not reachable from \texttt{retry}. The
structural baseline correctly excludes it but does not order the
remaining candidates.

\textbf{Table 2. Recovery candidates and conservative historical
support.}

\begin{longtable}[]{@{}
  >{\raggedright\arraybackslash}p{(\columnwidth - 10\tabcolsep) * \real{0.1304}}
  >{\raggedleft\arraybackslash}p{(\columnwidth - 10\tabcolsep) * \real{0.1739}}
  >{\raggedleft\arraybackslash}p{(\columnwidth - 10\tabcolsep) * \real{0.1739}}
  >{\raggedleft\arraybackslash}p{(\columnwidth - 10\tabcolsep) * \real{0.1739}}
  >{\raggedleft\arraybackslash}p{(\columnwidth - 10\tabcolsep) * \real{0.1739}}
  >{\raggedleft\arraybackslash}p{(\columnwidth - 10\tabcolsep) * \real{0.1739}}@{}}
\toprule\noalign{}
\begin{minipage}[b]{\linewidth}\raggedright
Candidate
\end{minipage} & \begin{minipage}[b]{\linewidth}\raggedleft
Raw success
\end{minipage} & \begin{minipage}[b]{\linewidth}\raggedleft
n
\end{minipage} & \begin{minipage}[b]{\linewidth}\raggedleft
Wilson lower 95\%
\end{minipage} & \begin{minipage}[b]{\linewidth}\raggedleft
Reachable from \texttt{retry}
\end{minipage} & \begin{minipage}[b]{\linewidth}\raggedleft
CT-RAG rank
\end{minipage} \\
\midrule\noalign{}
\endhead
\bottomrule\noalign{}
\endlastfoot
ProviderFallback & 0.80 & 10 & \textbf{0.490} & Yes & \textbf{1} \\
ScriptPatch & 0.20 & 10 & 0.057 & Yes & 2 \\
CacheReset & 0.25 & 4 & 0.046 & Yes & 3 \\
ManualPatch & 0.10 & 10 & 0.018 & Yes & 4 \\
GlobalFix & \textbf{0.90} & 10 & \textbf{0.596} & \textbf{No} & excluded \\
\end{longtable}

The order of \texttt{ScriptPatch} and \texttt{CacheReset} illustrates
why empirical point rate alone is not the recovery score: Wilson support
for 2/10 (0.057) exceeds 1/4 (0.046), despite the latter's higher raw
point estimate. \texttt{ProviderFallback} remains clearly separated:
its lower bound is 0.490, a gap of 0.433 from \texttt{ScriptPatch} and
0.444 from \texttt{CacheReset}.

\textbf{Table 3. Recovery comparison with metric scope made explicit.}

\begin{longtable}[]{@{}llrr@{}}
\toprule\noalign{}
Method & Output & Precision@1 & Precision@4 \\
\midrule\noalign{}
\endhead
\bottomrule\noalign{}
\endlastfoot
Terrain-only & HumanIntervention & 0.0 & --- \\
Recency-only & ManualPatch & 0.0 & --- \\
Raw success, global & GlobalFix & 0.0 & --- \\
Structural reachable ceiling & unranked set of four & n/a & 0.25 \\
CT-RAG \texttt{RECOVERY} & \textbf{ProviderFallback} & \textbf{1.0} & 0.25 \\
\end{longtable}

The top-1 result alone is insufficient for a healer that may consume
multiple suggestions. The measured false-rescue curve is therefore
reported directly:

{[} FR@1=0.00,\quad FR@2=0.50,\quad FR@3=0.667,\quad FR@4=0.75. {]}

This degradation is not hidden or thresholded away. It shows that the
current mechanism is strong at selecting the first candidate in this
fixture but does not solve the broader problem of returning a clean
multi-candidate set. A production healer should therefore either consume
top-1 under stricter confidence policy or add further constraints before
acting on lower-ranked suggestions.

The central inference from this experiment is again narrow:
\textbf{global historical success is insufficient when the
highest-success route is not reachable from the current state.} The
causal topology supplies an eligibility relation not present in terrain,
recency, or success statistics.

\subsubsection{Reachability-gate isolation
check}\label{reachability-gate-isolation-check}

After the numerical evidence freeze used for Tables 1--3, we executed
the preregistered mutation that changes only feasibility: an edge
\texttt{retry\ -\textgreater{}\ global\_fix} is added while the
terminal success statistics remain 9/10 for \texttt{GlobalFix} and 8/10
for \texttt{ProviderFallback}. Under this mutation, both candidates are
causally reachable and \texttt{GlobalFix} becomes the top-ranked
recovery, with its Wilson lower bound (0.596) exceeding
\texttt{ProviderFallback} (0.490). The test passed in GitHub Actions CI
\textbf{\#213}, head commit
\textbf{\texttt{9ed7878813b951ca04428e27868d5a4427e8011f}}, on Python
3.11, 3.12, and 3.13.

This isolation check is intentionally reported separately from the CI
\#210 numerical freeze: it adds no new headline effect size. Its role is
falsification-oriented. The observed rank flip supports the
interpretation that causal reachability acts as an eligibility gate
rather than as an undeclared continuous penalty hidden inside the
recovery score.

\subsection{8.3 Secondary result: matched-FPR early
degradation}\label{secondary-result-matched-fpr-early-degradation}

The original fixed infrastructure threshold produced an apparent +2-day
lead at the default behavioral configuration. That comparison was
demoted because it compared unlike detector operating rules. The
stronger ablation applies the \textbf{same sustained-change detector
logic} to behavioral and infrastructure signals and calibrates the
infrastructure threshold on the same null generator. At
(\theta=0.10) with two sustained windows, behavioral change is detected
on day 8 and infrastructure change on day 9, yielding a \textbf{+1-day
relative lead}.

The current null generator, however, yields FPR = 0 for both detectors
throughout the tested grid. Consequently, the observed FPR gap of 0 is
\textbf{degenerate}: it does not demonstrate informative matching at a
non-zero operating point. The protection of the +1-day result comes from
comparing the same detector logic and from its robustness across
parameter/phase perturbations, not from evidence that non-zero
false-positive rates have been successfully matched.

Across the full 28-cell grid:

\begin{itemize}
\tightlist
\item
  16 cells have positive lead;
\item
  achieved FPR is 0 for both detectors in all 28 cells, so the FPR gap
  is 0 but non-informative as an operating-point match;
\item
  the 16 positive cells form one contiguous region in the tested
  parameter lattice;
\item
  median lead across the 28 cells is +1 day;
\item
  the lead distribution ranges from -1 to +3 days.
\end{itemize}

The positive region is:

{[}

\begin{aligned}
&\theta=0.05, &&s=1..4\\
&\theta=0.075,&&s=1..3\\
&\theta=0.10, &&s=1..3\\
&\theta=0.125,&&s=1..3\\
&\theta\in\{0.15,0.175,0.20\},&&s=1.
\end{aligned}

{]}

This shape matters more than the fraction 16/28: positive lead persists
under stricter sustained-window requirements only at lower drift
thresholds, indicating a modest signal that requires sensitivity rather
than a large, obvious separation.

The phase-offset ablation produces paired leads
\texttt{{[}1,\ 1,\ 1,\ 1{]}} for phase fractions 0, 0.25, 0.50, and
0.75. At 0.50 and 0.75, both detectors move one day later, while their
difference remains +1. We therefore do \textbf{not} claim an absolute
detection date such as ``day 8'' as phase invariant. The stable quantity
is the relative lead under this synthetic daily-resolution fixture.

These observations should not be described as production forecasting.
They show an early-warning mechanism under a controlled generative
pattern, not predictive validity on independently sampled incidents.

\section{9. Mechanism Validation}\label{mechanism-validation}

Property tests help distinguish design invariants from performance
claims.

\subsubsection{Obsolete-path decay}\label{obsolete-path-decay}

A historically reinforced legacy recovery edge peaks at influence 5.5
and decays to approximately 0.309 while its historical transition count
remains 10 and the edge remains present. The replacement async path
reaches influence approximately 3.285 and its recovered terminal ranks
above the legacy recovered terminal. This validates the separation
between navigational erosion and historical deletion.

\subsubsection{Basin attraction drift}\label{basin-attraction-drift}

After a recurrent compensation cycle is introduced and repeatedly
observed, internal SCC observations increase from 0 to 80 and discovered
attractor confidence increases from 0.5 to 1.0. The affected-probe
membership gain for the problematic basin is 0.75. These are
deterministic fixture-level mechanism checks; they are not estimates of
real-world basin dynamics.

\subsubsection{Out-of-order
reconciliation}\label{out-of-order-reconciliation}

Before the delayed parent arrives, the child has no incoming causal
edge. After the parent arrives, exactly one event-provenance causal edge
is reconciled, despite the 336-hour event-time gap. This invariant
underlies the comparative RQ1 result.

\section{10. Threats to Validity}\label{threats-to-validity}

\subsection{10.1 Internal validity}\label{internal-validity}

The primary experiments are adversarial but hand-constructed. Their
value is mechanistic isolation: each baseline is given a signal that
should make it competitive, and the fixture tests whether causal
topology contributes distinct information. Hand construction also
creates a risk that the score function and fixture co-evolve. The
\texttt{GlobalFix} reachability mutation described above is therefore
important because it can falsify an undeclared scoring dependency.

The recovery score contains a fixed coefficient (\lambda=0.35). We do
not claim this value is optimal. The top-1 ordering in the current
fixture is robust to the large Wilson separation between
\texttt{ProviderFallback} and low-quality reachable distractors, but
broader sensitivity analysis remains appropriate.

\subsection{10.2 Construct validity}\label{construct-validity}

\texttt{causation\_id} is treated as authoritative runtime provenance.
This is valid only to the extent that the emitting runtime maintains the
field correctly. CT-RAG guarantees provenance-preserving retrieval, not
the truth of incorrectly emitted provenance.

Likewise, ``basin'' is an engineered graph construct based on sinks,
recurrent SCCs, and bounded reachability. It is not a claim about
dynamical systems in the continuous mathematical sense.

The Wilson lower bound is a conservative support score, not a posterior
probability and not a causal success probability under intervention.

\subsection{10.3 External validity}\label{external-validity}

The strongest results are controlled synthetic fixtures. They establish
that runtime topology can encode useful information under specific
adversarial geometries, but they do not establish improvements on
production incidents, LongMemEval, AIOps traces, or arbitrary agent
tasks.

The repository contains adapters and external-telemetry scaffolding, but
a versioned external microservice fixture sufficient for a
publication-level external-validity claim is not yet part of the
evidence reported here. In telemetry without authoritative causation
metadata, CT-RAG deliberately operates with weaker temporal/behavioral
relations and should not be described as recovering authoritative
causality.

\subsection{10.4 Statistical validity}\label{statistical-validity}

The recovery and causal-gap experiments are deterministic mechanism
fixtures, not sampled populations; confidence intervals over repeated
identical runs would be meaningless. Their evidence comes from
adversarial counter-baselines and falsifiable mutation tests rather than